\section{Evidence-Bound Research State Machine}
\label{sec:method}

\subsection{Campaign state and counted successors}

A campaign is a persistent state machine rather than a sequence of unrelated
prompts. Figure~\ref{fig:campaign_pipeline} shows the ten implemented states.
The campaign manifest records the policy, resource limits, current state,
round index, intervention count, and lineage hash. Each round manifest stores
the parent-result hash and content hashes for its structured artifacts.

\begin{figure*}[t]
  \centering
  \includegraphics[width=.98\textwidth]{figures/campaign-pipeline.pdf}
  \Description{Ten stages from observation through report. A feedback arrow
  connects a negative report to a new diagnosis in a later round.}
  \caption{Persistent campaign state machine. Outcome data become available
  only after the freeze. A negative decision can start a later diagnosis but
  cannot modify the adjudicated round.}
  \label{fig:campaign_pipeline}
\end{figure*}

Let \(r_i\) denote round \(i\), \(e_{i-1}\) its parent result,
\(h(\cdot)\) a SHA-256 content hash, and \(m_i\) the proposed mechanism
family. The state machine counts \(r_i\) as a scientific successor only when five conditions
hold:

\begin{enumerate}
  \item the proposal records \(h(e_{i-1})\) and a typed diagnosis of the
  parent failure;
  \item \(m_i\) changes a stated causal or statistical operation, rather than
  only parameters or wording;
  \item code, data identifiers, idempotency identifiers, metrics, stopping
  rules, and the adjudicator are frozen before outcome access;
  \item execution creates new observations and a deterministic decision; and
  \item both the research report and any failure report remain in the lineage.
\end{enumerate}

An identical rerun is reproduction, a manuscript edit is reporting, and a
prompt edit without new evidence is neither. The counted object is therefore
a parent-linked transition through a new empirical contract.

\subsection{Evidence graph and adjudication boundary}

The evidence graph contains typed claim, artifact, activity, decision,
limitation, and contradiction records. A report claim must resolve to a file,
JSON pointer or line locator, content hash, and allowed scope. Claims without
a live edge fail the deterministic report gate. This local graph follows the
same entity-activity-agent separation used by PROV-O
~\cite{lebo2013provo}, but the current serialization is project-specific.

The local language model receives the parent diagnosis, permitted development
context, mechanism constraints, and an output schema. In this study the model
is from the Qwen family~\cite{yang2025qwen3}. It may propose a diagnosis,
mechanism, falsifier, or concise policy. It cannot choose unseen data after
outcome access, calculate the acceptance statistic, alter a threshold, or
mark a contribution as accepted. A timeout or malformed response selects a
recorded deterministic fallback and cannot change numerical adjudication.

All decisions are functions of frozen records. The mechanism screen checks
that a successor changes the declared mechanism family. The development gate
calculates the registered statistic and validates cell coverage. The freeze
records hashes of candidate code, data identities, parameter space, metrics,
stopping rules, and adjudicator source. The unseen evaluator writes one result
for that identity. The adjudicator maps it to positive, negative, or invalid.

For preregistration \(P_i\), adjudicator source \(A_i\), and parent result
\(e_{i-1}\), the freeze identity is
\[
F_i=h\!\left(h(P_i)\,\|\,h(A_i)\,\|\,h(e_{i-1})\right).
\]
The result stores \(F_i\), and the decision stores the result hash. This chain
detects later substitution. It does not determine novelty, experimental
adequacy, or truth.

Figure~\ref{fig:evidence_boundary} separates the three statistical and
governance layers that were conflated in the original internal report.

\begin{figure*}[t]
  \centering
  \includegraphics[width=.98\textwidth]{figures/evidence-boundary-map.pdf}
  \Description{Three columns separate the historical 30-cell internal gate,
  the additive ten-task analysis, and a future independent confirmation. The
  first two are complete and the third remains required.}
  \caption{Evidence boundary. Deterministic repetitions can test stability,
  but only the ten task outcomes enter the additive uncertainty analysis.
  Neither layer replaces an independently authored confirmation and a human
  publication decision.}
  \label{fig:evidence_boundary}
\end{figure*}

\subsection{Permissioned project memory}

The \vault{} is an Obsidian-compatible and Git-friendly project memory. It
stores literature notes, topic indexes, experiment records, failure cases,
skill and strategy cards, claim links, and version history. Retrieval is
state-dependent. A failure-analysis card may inform a later diagnosis, while
an unseen metric cannot enter the proposal context for the current round.
Every retrieved memory item has a stable path and is recorded in the round
manifest.

This memory design differs from an unconstrained conversation history.
Negative mechanisms remain addressable, reports link back to their runs and
decisions, and permission labels define which state may consume an artifact.
The controlled ablation removes this memory channel for tasks whose repair
depends on a prior evidence card.

\subsection{Terminal artifacts and human controls}

Each completed round exports its hypothesis, preregistration, experiment
manifest, raw metrics, validation record, failure analysis, research report,
loop report, claim-evidence map, figures, decision, and manuscript version.
The package also records tool versions, build commands, hashes, and source
status. These objects support inspection and rerunning.

Five decisions remain outside the state machine: target venue, authorship,
license resolution, AI-use disclosure, and external distribution or
submission. The package stores each as unresolved or false. A successful
internal build cannot change those values, and self-iteration cannot modify
their policy.
